\documentclass[12pt,a4paper]{article}

% --- Paquetes Esenciales ---
\usepackage[utf8]{inputenc}
\usepackage[T1]{fontenc}
\usepackage[spanish]{babel}
\usepackage{amsmath}
\usepackage{amssymb}
\usepackage{geometry}
\usepackage{graphicx}
\usepackage{booktabs}
\usepackage{multirow}
\usepackage[hidelinks]{hyperref}

% --- Configuración de la página ---
\geometry{a4paper, margin=2.5cm}

\setcounter{secnumdepth}{0} % Desactiva la numeración automática de secciones

\begin{document}

\begin{center}
    \textbf{\Large UDO – NÚCLEO DE SUCRE}\\
    \vspace{0.2cm}
    \textbf{\large ESCUELA DE CIENCIAS – DPTO. DE FÍSICA}\\
    \vspace{0.2cm}
    \textbf{\large LABORATORIO I DE FÍSICA}
\end{center}

\vspace{0.5cm}
\hrule
\vspace{0.5cm}

\tableofcontents
\newpage

\begin{center}
    \textbf{\LARGE LA MEDICIÓN Y EL ERROR}
\end{center}
\vspace{0.3cm}

La física es una ciencia y como tal está fundamentada en la observación y la medición. Guiado por la curiosidad, el físico observa al mundo que lo rodea, percibe los hechos, mide o evalúa y luego saca conclusiones para elaborar las leyes de la física.

La medición nunca es exacta. Las magnitudes físicas pueden ser determinadas experimentalmente mediante medidas o combinaciones de medidas. Estas medidas tienen una incertidumbre derivada de las características de los aparatos que son utilizados en su determinación.

En el laboratorio de física no vamos a elaborar nuevas leyes sino simplemente vamos a comprobar el cumplimiento de algunas de ellas. Por ello, es fundamental que el estudiante aprenda a medir, conozca los orígenes de la incertidumbre de sus medidas y sepa expresar numéricamente los resultados obtenidos.

Toda medida de una cantidad física viene siempre acompañada de una unidad dependiendo del sistema en el cual se mide. Por ejemplo, una longitud puede medirse en metros, pulgadas o en años luz. Cada sistema de medida tiene un patrón. El trabajo en el laboratorio implica medir magnitudes físicas mediante la utilización de instrumentos calibrados contra un patrón de medidas. Este patrón de medida debe ser universal. Si cada persona tuviera su propio patrón no podrían establecerse comparaciones entre medidas realizadas por diferentes personas, a menos que se conocieran las equivalencias entre cada uno de los patrones involucrados.

Es bueno destacar que como ningún instrumento es perfecto y que tampoco los humanos lo son, toda medida tiene asociada una incertidumbre. Obsérvese que, debido a lo antes expuesto, si una cantidad física se mide varias veces, por una misma persona, no siempre se obtiene el mismo resultado. Es decir, que el resultado de una medición experimental es siempre un valor aproximado. Por lo tanto, al determinar una magnitud física no se encuentra su valor exacto sino el valor más probable de la misma. Este valor se encuentra contenido en un intervalo acotado dentro del cual están todas las posibles medidas que alguna persona pueda realizar.

A menudo las palabras error e incerteza son usadas como sinónimos para denotar la incertidumbre en la medición, sin embargo, estas dos palabras tienen significados diferentes en el contexto de medidas. Cuando realizamos una medida, tratamos de determinar el valor verdadero de la cantidad de interés, pero el resultado que tenemos difiere del valor ``verdadero''. Esta diferencia entre el valor verdadero y el resultado de nuestra medida es el error en la medida. Por otra parte, al repetir las medidas éstas no necesariamente son idénticas por lo que existe siempre una incerteza en nuestros resultados.

Dentro de los objetivos del laboratorio de física están: enseñar a usar los instrumentos de medidas, entrenar al estudiante para que reduzca al mínimo los errores humanos en el proceso de medición y, principalmente, adiestrar al futuro profesional para que sepa determinar y expresar las incertezas de su medición. La evaluación del laboratorio de física hace énfasis en la capacidad de un estudiante en dar estimación del tamaño del error de una medida individual.

\vspace{0.5cm}
\hrule
\vspace{0.5cm}

\section{I.- LA APRECIACIÓN DE UN INSTRUMENTO}

Se define como \textbf{apreciación de un instrumento} a la medida más pequeña que se puede realizar con dicho equipo. Alternativamente pudiera decirse que la apreciación es la menor diferencia que puede obtenerse entre valores de una magnitud, medidos con dicho instrumento. Por lo general los instrumentos vienen calibrados en divisiones y la apreciación es la menor división de la escala del instrumento. Por ejemplo, una regla cuya menor división es 1 mm tiene una apreciación de 1 mm.

Cuando se obtiene una medida mediante un instrumento que tiene una escala calibrada en divisiones, se aproxima la lectura a la división más cercana. Por esto, el máximo error que se puede cometer en dicha medición es de más o menos la apreciación del instrumento. Debido a ello toda medida presenta siempre un error determinado por la precisión del instrumento.

En los instrumentos que tienen una sola escala, como la regla, la apreciación se determina de la siguiente manera: se escogen dos valores sobre la escala, que pueden ser consecutivos o no, y se calcula la apreciación mediante la fórmula:

\begin{equation}
\text{Apreciación } (A) = \frac{\text{Lectura superior} - \text{Lectura inferior}}{\text{N}^\circ \text{ de divisiones entre las lecturas}} = \frac{L_s - L_i}{n} \tag{1}
\end{equation}

\vspace{0.5cm}
\hrule
\vspace{0.5cm}

\section{II.- PRECISIÓN Y EXACTITUD}

Además de la apreciación, existen otras dos palabras relacionadas con las medidas que obtenemos mediante un instrumento: \textit{precisión} y \textit{exactitud}. La precisión se refiere a la cercanía de los valores medidos entre sí. La exactitud se refiere a la cercanía de las mediciones al valor verdadero de una magnitud física; la exactitud está asociada con la calibración del instrumento que utilizamos para medir. Veamos lo que se acaba de explicar con un ejemplo:

\begin{quote}
``Se desea medir el tiempo con un cronómetro de apreciación \textbf{0,01 s} pero que se adelanta un cierto valor cada hora. También se cuenta con un reloj de pulsera con apreciación de \textbf{1 s} que no se adelanta. En este caso podemos decir que aun siendo el cronómetro más preciso que el reloj de pulsera sus resultados son menos exactos. Es por ello que es deseable que la calibración del instrumento sea tan buena como la apreciación del mismo''.
\end{quote}

En las figuras a, b y c se muestra de manera esquemática los conceptos de precisión y exactitud:
\begin{itemize}
    \item En la \textbf{figura a}, las medidas obtenidas son más precisas pero menos exactas.
    \item En la \textbf{figura b}, las medidas son más exactas pero menos precisas.
    \item La \textbf{figura c} expone medidas precisas y exactas.
\end{itemize}

La precisión de un instrumento es inversamente proporcional a la apreciación de un instrumento:

\begin{equation}
P = \frac{1}{A} = \frac{n}{L_s - L_i} \tag{2}
\end{equation}

\vspace{0.5cm}
\hrule
\vspace{0.5cm}

\section{III.- EXPRESIÓN DE UNA LECTURA}

Supongamos que se trata de medir una cantidad cuyo valor verdadero es $X$. Al realizarse la medida el observador obtiene la lectura $X$ con un error $\Delta X$. Esto significa, que el aparato de medida informa que el valor de la cantidad medida pertenece al intervalo de confianza:

\begin{equation}
(X - \Delta X) \leq X \leq (X + \Delta X) \tag{3}
\end{equation}

El resultado de cualquier magnitud medida debe escribirse en la forma:

\begin{equation}
(X \pm \Delta X) \tag{4}
\end{equation}

donde $X$ representa el \textit{valor de la medida} y $\Delta X$ es el \textit{error en la medida}.

Al escribir una medida también debe tenerse en cuenta la potencia que la afecta y su unidad; tanto la potencia como la unidad se escriben fuera del paréntesis. Por ejemplo:

\[ (10,0 \pm 0,2) \times 10^3\ \text{kg} \]

\textit{Obsérvese que el número de decimales de la cantidad $X$ deber ser igual al número de decimales del error $\Delta X$.}

\subsection*{Significado del número cero cuando este se usa como decimal}

Según la aritmética resulta lo mismo escribir el número nueve como $9$ o como $9,0$ o sea $9 = 9,0$. En el caso de las mediciones que se hacen en el laboratorio, estas dos expresiones no significan lo mismo por lo que no es igual escribir, por ejemplo, $9\ \text{mm}$ que $9,0\ \text{mm}$. Se dice que una de ellas es más precisa que la otra; es decir, ésta afectada por un error menor. Por ejemplo:

Supongamos que escribimos una medida como: 
\[ (9 \pm 1) \times 10^1\ \text{mm} \]

En este caso la medida son noventa milímetros más o menos diez milímetros, o sea que los valores máximos y mínimos de esa medida $X$ están comprendidos entre:

\[ (X + \Delta X) = (90 + 10)\ \text{mm} \]
\[ (X - \Delta X) = (90 - 10)\ \text{mm} \]
\[ 80\ \text{mm} \leq X \leq 100\ \text{mm} \]

La incertidumbre entre el valor máximo y el mínimo es de $20\ \text{mm}$.

Supongamos ahora, que escribimos la medida como: 
\[ (9,0 \pm 0,1) \times 10^1\ \text{mm} \]

Los valores máximos y mínimos de esa medida $X$ están comprendidos entre:

\[ (X + \Delta X) = (90 + 1)\ \text{mm} \]
\[ (X - \Delta X) = (90 - 1)\ \text{mm} \]
\[ 89\ \text{mm} \leq X \leq 91\ \text{mm} \]

En este caso, la incertidumbre es de $2\ \text{mm}$.

Puede observarse que la segunda medida es más precisa que la primera, aun cuando ambas reportan el mismo valor. Se dice entonces que el error hace la diferencia en cuanto a la precisión de la medida.

\vspace{0.5cm}
\hrule
\vspace{0.5cm}

\section{IV.- ERRORES EXPERIMENTALES}

Cuando se realizan medidas de una magnitud nos encontramos con dos tipos de errores: los \textbf{sistemáticos} y los \textbf{aleatorios}. Los \textit{errores sistemáticos} pueden reducirse al mínimo controlando la experimentación, pero los \textit{errores aleatorios} están fuera del control de la persona que mide.

Los \textbf{errores sistemáticos} son errores que producen un resultado que difiere del valor verdadero en una cantidad fija. Estos son debidos a defectos en el equipo de medida o al uso incorrecto de las técnicas de medición. Los errores sistemáticos pueden ser identificados y corregidos.

Como ejemplos de errores sistemáticos, podemos mencionar:
\begin{enumerate}
    \item \textbf{Errores de calibración de los instrumentos.} Por ejemplo:
    \begin{itemize}
        \item[a)] Posición incorrecta de la aguja de un instrumento analógico o mecánico.
        \item[b)] Falta de calibración interna de los aparatos electrónicos.
        \item[c)] Mala calibración por construcción defectuosa.
    \end{itemize}
    \item \textbf{Mala escogencia del rango de un instrumento para medir una magnitud.} Cuando un instrumento permite medir una magnitud en diferentes rangos se debe tener especial cuidado, ya que cuando se pasa de un rango a otro los errores no son los mismos. Siempre se sugiere usar el rango más pequeño posible que permita medir la magnitud, con el fin de minimizar el error asociado a la medida.
    \item \textbf{Errores de observación.} La mala ubicación del observador frente al instrumento de medida representa un error sistemático (error de paralaje).
\end{enumerate}

Los \textbf{errores aleatorios o al azar} son un tipo de error consecuencia de la naturaleza de la magnitud que se desea medir y de las influencias que pueda tener el medio sobre la magnitud. Corresponden a este tipo de error:

\begin{enumerate}
    \item \textbf{Errores de apreciación.} Por ejemplo, en la estimación de la fracción de la menor división de la escala. 
    \item \textbf{Errores debido a condiciones que fluctúan.} Por ejemplo, variaciones en la red de energía eléctrica, en la temperatura, etc.
    \item \textbf{Indefiniciones físicas.} Aunque el proceso de medida fuera perfecto, la repetición de una medida puede dar valores diferentes, porque la magnitud a medirse, no está del todo definida (ej. el diámetro de una esfera que no es perfectamente lisa).
\end{enumerate}

\vspace{0.5cm}
\hrule
\vspace{0.5cm}

\section{V.- DEFINICIONES}

Si hacemos varias mediciones, el mejor estimado del valor ``verdadero'' es el promedio de todas las medidas. Sea $X$ la cantidad que queremos medir y supongamos que realizamos ``n'' medidas de esta cantidad. El \textbf{valor promedio} de $X$, el cual denotamos por $\bar{X}$ está dada por la ecuación:

\begin{equation}
\bar{X} = \frac{\sum_{i=1}^{n} X_i}{n} \tag{5}
\end{equation}

donde $X_i$ representa la i-ésima medida de la cantidad $X$.

¿Cómo determinar el error de esta medida promedio? Esta cantidad, denotada por $\Delta X$, se llama \textbf{desviación de una medida} y está dada por la ecuación:

\begin{equation}
\Delta X_i = X_i - \bar{X} \tag{6}
\end{equation}

Una manera de corregir que las desviaciones se cancelen es promediar el valor absoluto de las desviaciones. Este promedio se llama \textbf{desviación absoluta media} ($\overline{\Delta X}$):

\begin{equation}
\overline{\Delta X} = \frac{\sum_{i=1}^{n} | X_i - \bar{X} |}{n} \tag{7}
\end{equation}

En los casos de muchas medidas se utiliza un parámetro estadístico llamado \textbf{desviación típica del valor medio} ($\sigma_m$) y se define como:

\begin{equation}
\sigma_m = \sqrt{ \frac{\sum_{i=1}^{n} (X_i - \bar{X})^2}{n(n-1)} } \tag{8}
\end{equation}

A $\sigma_m$ también se le llama \textit{error cuadrático medio del promedio}. En el laboratorio el resultado se escribirá como:

\begin{equation}
(\bar{X} \pm 3\sigma_m) \tag{9}
\end{equation}

\textbf{Desviación media relativa de una (serie de) medida(s) o error relativo ($E_r$)}

Se define como la desviación dividida entre el valor medio:

\begin{equation}
E_r = \frac{\Delta X}{\bar{X}} \quad \text{o} \quad E_r = \frac{\overline{\Delta X}}{\bar{X}} \tag{10}
\end{equation}

\textbf{Desviación media porcentual de una (serie de) medida(s) o error porcentual ($E_{\%}$)}

Se define como el error relativo multiplicado por 100:

\begin{equation}
E_{\%} = \frac{\Delta X}{\bar{X}} \times 100 \quad \text{o} \quad E_{\%} = \frac{\overline{\Delta X}}{\bar{X}} \times 100 \tag{11}
\end{equation}

\subsection*{Ejemplo:}
Supongamos que deseamos medir el diámetro de un cilindro.

\begin{table}[h!]
\centering
\begin{tabular}{|c|c|c|}
\hline
\textbf{N$^\circ$} & \textbf{(d $\pm$ 0,01) mm} & \textbf{$| d_i - \bar{d} |$} \\ \hline
1 & 9,93 & 0,00 \\ \hline
2 & 9,90 & 0,03 \\ \hline
3 & 9,94 & 0,01 \\ \hline
4 & 9,95 & 0,02 \\ \hline
5 & 9,91 & 0,02 \\ \hline
6 & 9,93 & 0,00 \\ \hline
7 & 9,92 & 0,01 \\ \hline
8 & 9,90 & 0,03 \\ \hline
9 & 9,95 & 0,02 \\ \hline
10 & 9,92 & 0,01 \\ \hline
\textbf{Prom.} & \textbf{9,93} & \textbf{0,02} \\ \hline
& $\bar{d}$ & $\overline{\Delta d}$ \\ \hline
\end{tabular}
\end{table}

\textbf{Diámetro = (9,93 $\pm$ 0,02) mm}

\vspace{0.5cm}
\hrule
\vspace{0.5cm}

\section{VI.- Número de medidas}

\begin{enumerate}
    \item Si se repite la medida una o varias veces y siempre resulta el mismo valor, entonces se especifica un solo dato. Para un solo dato el error es la apreciación del instrumento.
    \item Si se mide una segunda vez y da un dato extremo (máximo o mínimo), se calcula el valor promedio $\bar{X}$ y se toma como error $\Delta X_m$:
    \begin{equation}
    \Delta X_m = \frac{X_{max} - X_{min}}{2} \tag{12}
    \end{equation}
    El resultado final se escribe:
    \begin{equation}
    X = (\bar{X} \pm \Delta X_m) \tag{13}
    \end{equation}
    \item Cuando la magnitud fluctúa de manera aleatoria, se aplica tratamiento estadístico.
    \begin{itemize}
        \item \textbf{A.} Para cualquier número de datos se puede aplicar la representación:
        \begin{equation}
        X = (\bar{X} \pm \overline{\Delta X}) \tag{14}
        \end{equation}
        \item \textbf{B.} Tratamiento riguroso (Desviación Típica):
        \begin{enumerate}
            \item Para $n > 20$ datos, el resultado está dado por la relación (9).
            \item Para $2 < n \leq 20$ datos, se recomienda el procedimiento (A).
        \end{enumerate}
    \end{itemize}
\end{enumerate}

\vspace{0.5cm}
\hrule
\vspace{0.5cm}

\section{VII.- PROPAGACIÓN DE ERRORES}

En la práctica, muchas magnitudes se calculan empleando una fórmula conocida $F = f(x, y, z, \dots)$ \textbf{(15)}.
La diferencial de la función $F$ es:

\begin{equation}
dF = \frac{\partial F}{\partial x} dx + \frac{\partial F}{\partial y} dy + \frac{\partial F}{\partial z} dz + \dots \tag{16}
\end{equation}

Tomando los diferenciales como desviaciones ($\Delta x \approx dx$):

\begin{equation}
\Delta F = \frac{\partial F}{\partial x} \Delta x + \frac{\partial F}{\partial y} \Delta y + \frac{\partial F}{\partial z} \Delta z + \dots \tag{17}
\end{equation}

Para el caso (A), el error máximo sobre $F$ estará dado por:

\begin{equation}
\Delta F = \left| \frac{\partial F}{\partial x} \right| \Delta x + \left| \frac{\partial F}{\partial y} \right| \Delta y + \left| \frac{\partial F}{\partial z} \right| \Delta z + \dots \tag{18}
\end{equation}

En el caso (B) (errores cuadráticos), el error está dado por:

\begin{equation}
\sigma_{mF} = \sqrt{ \left( \frac{\partial F}{\partial x} \right)^2 \sigma_{mx}^2 + \left( \frac{\partial F}{\partial y} \right)^2 \sigma_{my}^2 + \left( \frac{\partial F}{\partial z} \right)^2 \sigma_{mz}^2 + \dots } \tag{19}
\end{equation}

Si $x$ se mide 1 vez, y $y$ se mide $n > 20$ veces:

\begin{equation}
\Delta F = \left| \frac{\partial F}{\partial x} \right| \Delta x + \left| \frac{\partial F}{\partial y} \right| 3\sigma_{my} \tag{20}
\end{equation}

\subsection*{Ejemplos de Propagación:}
\begin{enumerate}
    \item Si $F = x + y$:
    \[ \Delta F = \left| \frac{\partial (x + y)}{\partial x} \right| \Delta x + \left| \frac{\partial (x + y)}{\partial y} \right| \Delta y = \Delta x + \Delta y \]
    \item Volumen de un cubo $V = a^3$, con $a = (42,56 \pm 0,01)\ \text{mm}$:
    \[ V = (42,56\ \text{mm})^3 = 77,09 \times 10^3\ \text{mm}^3 \]
    \[ \Delta V = \left| 3a^2 \right| \Delta a = 3 \cdot (42,56\ \text{mm})^2 \cdot 0,01\ \text{mm} = 0,05 \times 10^3\ \text{mm}^3 \]
    Volumen: $V = (77,09 \pm 0,05) \times 10^3\ \text{mm}^3$
    \item Densidad $\rho = \frac{m}{\pi \cdot r^2 \cdot L}$:
    \[ \Delta \rho = \left| \frac{1}{\pi \cdot r^2 \cdot L} \right| \Delta m + \left| - \frac{2m}{\pi \cdot r^3 \cdot L} \right| \Delta r + \left| - \frac{m}{\pi \cdot r^2 \cdot L^2} \right| \Delta L \]
\end{enumerate}

\subsection*{Actividades:}
\begin{enumerate}
    \item Dadas las siguientes expresiones, $F(x, y)$, hallar la fórmula de la propagación del error, si se suponen conocidos los errores $\Delta x$ y $\Delta y$.
    
    \vspace{0.3cm}
    \begin{tabular}{@{}llll@{}}
        a) $F = x - y$ & \hspace{0.5cm} b) $F = xy$ & \hspace{0.5cm} c) $F = \frac{x}{y}$ & \hspace{0.5cm} d) $F = x^n$ \\[0.4cm]
        e) $F = \frac{x + y}{x - y}$ & \hspace{0.5cm} f) $F = \ln x$ & \hspace{0.5cm} g) $F = \frac{xy}{1 + xy}$ & \hspace{0.5cm} i) $F = \frac{\pi \cdot x^{2y}}{4}$
    \end{tabular}
    \vspace{0.3cm}
    
    \item Dada la función $F = 3xy^2z$, con:
    \begin{itemize}
        \item[] $x = (5,22 \pm 0,03)\ \text{m}$
        \item[] $y = (12,5 \pm 0,5)\ \text{m}$
        \item[] $z = (145,4 \pm 0,1)\ \text{m}$
    \end{itemize}
    Calcule el error relativo $\frac{\Delta F}{F}$.
    
    \item Midiendo la masa $m$ y el volumen $V$ de un sólido se encuentra:
    \begin{itemize}
        \item[] $m = (177,54 \pm 0,02)\ \text{g}$
        \item[] $V = (22,44 \pm 0,01)\ \text{cm}^3$
    \end{itemize}
    Calcule la densidad $\rho = \frac{m}{V}$, con el error.
\end{enumerate}

\subsection*{Soluciones a las Actividades:}
\begin{enumerate}
    \item \textbf{Fórmulas de propagación de error (basado en error máximo absoluto):}
    \begin{itemize}
        \item[a)] $\Delta F = \Delta x + \Delta y$
        \item[b)] $\Delta F = |y|\Delta x + |x|\Delta y$
        \item[c)] $\Delta F = \frac{\Delta x}{|y|} + \frac{|x|\Delta y}{y^2}$
        \item[d)] $\Delta F = |n x^{n-1}| \Delta x$
        \item[e)] $\Delta F = \frac{2|y|\Delta x + 2|x|\Delta y}{(x-y)^2}$
        \item[f)] $\Delta F = \frac{\Delta x}{|x|}$
        \item[g)] $\Delta F = \frac{|y|\Delta x + |x|\Delta y}{(1+xy)^2}$
        \item[i)] Asumiendo la fórmula del volumen de un cilindro $F = \frac{\pi x^2 y}{4}$: $\Delta F = \frac{\pi}{4}(2|xy|\Delta x + x^2\Delta y)$.
    \end{itemize}
    \item \textbf{Error relativo de $F = 3xy^2z$:} \\
    $\frac{\Delta F}{F} = \frac{\Delta x}{x} + 2\frac{\Delta y}{y} + \frac{\Delta z}{z} = \frac{0,03}{5,22} + 2\left(\frac{0,5}{12,5}\right) + \frac{0,1}{145,4} \approx 0,0057 + 0,0800 + 0,0007 = \mathbf{0,0864}$
    \item \textbf{Cálculo de la densidad $\rho$:} \\
    $\rho = \frac{m}{V} = \frac{177,54}{22,44} \approx 7,91176\ \text{g/cm}^3$ \\
    $\frac{\Delta \rho}{\rho} = \frac{\Delta m}{m} + \frac{\Delta V}{V} = \frac{0,02}{177,54} + \frac{0,01}{22,44} \approx 0,00056$ \\
    $\Delta \rho = 7,91176 \times 0,00056 \approx 0,0044\ \text{g/cm}^3$ \\
    Resultado final redondeado: $\mathbf{\rho = (7,912 \pm 0,004)\ \text{g/cm}^3}$
\end{enumerate}

\vspace{0.5cm}
\hrule
\vspace{0.5cm}

\section{VIII.- NOTACIÓN CIENTÍFICA (NC)}

La notación científica expresa un número como el producto de un número entre 1 y 9 por una potencia de diez. 
\begin{itemize}
    \item Exponente \textbf{positivo} para números grandes.
    \item Exponente \textbf{negativo} para números pequeños.
\end{itemize}

Ejemplos:
\begin{itemize}
    \item $1.000.000.000 = 1 \times 10^9$
    \item $0,000000001 = 1 \times 10^{-9}$
    \item $576 = 5,76 \times 10^2$
    \item $987.600 = 9,876 \times 10^5$
\end{itemize}

\textbf{Suma y resta en la NC:}\\
Al sumar o restar, los exponentes deben ser iguales.\\
Ejemplo: $5,678 \times 10^8 + 6,1234 \times 10^5 \rightarrow 5,678 \times 10^8 + 0,0061234 \times 10^8 = 5,6841234 \times 10^8$

\vspace{0.3cm}
\textbf{Multiplicación y división en la NC:}\\
Ejemplo: $(4,2 \times 10^{10}) \cdot (5,205 \times 10^2) = (4,2 \cdot 5,205) \times 10^{10+2} = 21,9 \times 10^{12} = 2,19 \times 10^{13}$

\subsection*{Prefijos del SI}

\begin{table}[h!]
\centering
\begin{tabular}{|l|l|c|c|l|}
\hline
\textbf{Múltiplos / Sub} & \textbf{Prefijo} & \textbf{Símbolo} & \textbf{Factor} & \textbf{Equivalencia decimal} \\ \hline
\multirow{10}{*}{Múltiplos} 
& yotta & Y & $10^{24}$ & 1 000 000 000 000 000 000 000 000 \\ \cline{2-5}
& zetta & Z & $10^{21}$ & 1 000 000 000 000 000 000 000 \\ \cline{2-5}
& exa & E & $10^{18}$ & 1 000 000 000 000 000 000 \\ \cline{2-5}
& peta & P & $10^{15}$ & 1 000 000 000 000 000 \\ \cline{2-5}
& tera & T & $10^{12}$ & 1 000 000 000 000 \\ \cline{2-5}
& giga & G & $10^{9}$ & 1 000 000 000 \\ \cline{2-5}
& mega & M & $10^{6}$ & 1 000 000 \\ \cline{2-5}
& kilo & k & $10^{3}$ & 1 000 \\ \cline{2-5}
& hecto & h & $10^{2}$ & 100 \\ \cline{2-5}
& deca & da & $10^{1}$ & 10 \\ \hline
\multirow{8}{*}{Submúltiplos} 
& deci & d & $10^{-1}$ & 0.1 \\ \cline{2-5}
& centi & c & $10^{-2}$ & 0.01 \\ \cline{2-5}
& mili & m & $10^{-3}$ & 0.001 \\ \cline{2-5}
& micro & $\mu$ & $10^{-6}$ & 0.000 001 \\ \cline{2-5}
& nano & n & $10^{-9}$ & 0.000 000 001 \\ \cline{2-5}
& pico & p & $10^{-12}$ & 0.000 000 000 001 \\ \cline{2-5}
& femto & f & $10^{-15}$ & 0.000 000 000 000 001 \\ \cline{2-5}
& atto & a & $10^{-18}$ & 0.000 000 000 000 000 001 \\ \hline
\end{tabular}
\end{table}

\vspace{0.5cm}
\hrule
\vspace{0.5cm}

\section{IX.- CIFRAS SIGNIFICATIVAS (CS)}

Las cifras significativas (CS) son todos los dígitos leídos en el instrumento, contados desde la izquierda a partir del primer dígito diferente de cero.

\textbf{Reglas:}
\begin{enumerate}
    \item Todos los dígitos distintos de cero son significativos.
    \item Cualquier cero entre dos dígitos significativos es siempre significativo.
    \item Los ceros ubicados al final son significativos.
    \item \textbf{En adición/sustracción:} El resultado tiene tantos decimales como el sumando con menos decimales.
    \item \textbf{En multiplicación/división:} El resultado tiene tantas cifras significativas como el valor con menos cifras significativas.
\end{enumerate}

\textbf{Para el Error $\Delta X$:}
\begin{itemize}
    \item \textbf{Regla 2 (del Error):} El error $\Delta X$ no debe tener más de dos cifras significativas.
    \item \textbf{Regla 3 (Precisión del dato):} La última cifra significativa del error y del valor de la medida deben estar en la misma posición decimal (Ej: $23,5 \pm 0,2$).
\end{itemize}

\vspace{0.5cm}
\hrule
\vspace{0.5cm}

\section{X.- REDONDEO DE CIFRAS}

\textbf{Regla 1:} Cuando el primero de los dígitos que se desea suprimir es menor que 5, el último dígito que se mantiene no se modifica. Si es mayor o igual a 5, se aumenta en una unidad la última cifra conservada.
\begin{itemize}
    \item Ej: Redondear $3,246\ \text{g}$ a dos CS $\rightarrow$ $3,2\ \text{g}$
    \item Ej: Redondear $1,237\ \text{m}$ a tres CS $\rightarrow$ $1,24\ \text{m}$
\end{itemize}

\textbf{Regla 2:} El redondeo no debe hacerse en forma sucesiva, sino siempre con respecto a la cifra original.
\begin{itemize}
    \item Ej: Redondear $7,249\ \text{mm}$ a dos CS $\rightarrow$ $7,2\ \text{mm}$ \textit{(Incorrecto hacerlo sucesivo: $7,249 \rightarrow 7,25 \rightarrow 7,3$)}
\end{itemize}

\vspace{0.5cm}
\hrule
\vspace{0.5cm}

\section{BIBLIOGRAFÍA}
\begin{enumerate}
    \item Davies, D. 1960. Métodos estadísticos aplicados a la investigación. Aguilar, Madrid.
    \item Young, H. 1962. Statistical treatment of experimental data. McGraw-Hill, New York.
    \item Phillip, R. 1969. Data reduction and error analysis for the physical sciences. McGraw-Hill, Bevington.
    \item Chatfield, Ch. 1975. Statistics for technology. Chapman \& Hall, London.
    \item ``Laboratorio I de Física'' 1ra. ed. 1991. Libro de Laboratorio. Departamento de Física. Universidad de Oriente – Venezuela.
    \item \url{https://www.ebsco.com/research-starters/mathematics/scientific-notation}
\end{enumerate}

\end{document}
\end{document}